% ------------------------------------------------------------------------------
% Anexo · Manual del desarrollador
% Fuentes: package.json (raíz y apps/web), README.md, AGENTS.md,
% apps/web/supabase/AGENTS.md, apps/web/src/config/*.config.ts, apps/web/.env*,
% packages/i18n/README.md, .github/workflows/ci.yml. Estado: 20/10/2026.
% Los bloques lstlisting solo contienen caracteres ASCII (listings no admite
% UTF-8 multibyte sin configuración adicional).
% ------------------------------------------------------------------------------
\label{cap:manual-desarrollador}

Este manual se dirige a quien vaya a evolucionar PymeKit o a construir sobre
él un SaaS nuevo. Explica cómo preparar el entorno, cómo está organizado el
código, cómo se trabaja con la base de datos, cómo se añade una tabla de
negocio que aparezca en la consola, qué se configura sin tocar el núcleo y cómo
se prueba y se despliega. Para la arquitectura y las decisiones que hay detrás
de cada pieza, véanse los capítulos de Diseño y de Construcción del Sistema.

\section{Introducción}
\label{sec:dg-introduccion}

PymeKit es un monorepo con una única aplicación web que contiene el área
pública, el área privada de los clientes, la consola de administración y la
API del CMS, más un conjunto de paquetes reutilizables. Todo comparte una sola
base de datos PostgreSQL gestionada con Supabase, cuyas políticas RLS son la
frontera de seguridad. Las directrices completas para personas y agentes de IA
están en el fichero \texttt{AGENTS.md} de la raíz y en los de cada paquete;
este manual resume lo imprescindible.

\section{Preparación del entorno de trabajo}
\label{sec:dg-entorno}

\subsection{Requisitos}

\begin{itemize}
  \item Node.js 20.10 o superior (se recomienda la versión 24).
  \item pnpm 11 (la versión exacta se fija en el campo
        \texttt{packageManager} del \texttt{package.json} raíz).
  \item Docker, que usa la instancia local de Supabase. El usuario debe poder
        ejecutar Docker sin \texttt{sudo}.
  \item Git y acceso al repositorio privado del proyecto.
  \item Opcional: Stripe CLI (se ejecuta en un contenedor) para recibir los
        \textit{webhooks} de pago en local.
\end{itemize}

\paragraph{Advertencia de seguridad.} Si se desarrolla en un servidor con IP
pública, los puertos que publica Docker pueden quedar expuestos aunque haya un
cortafuegos, porque Docker inserta sus propias reglas. Hay que filtrar en la
cadena \texttt{DOCKER-USER} o enlazar los servicios a \texttt{127.0.0.1} y
acceder mediante un túnel SSH (ADR-010, incidencia B-04).

\subsection{Arranque local}

\begin{lstlisting}[language=bash]
git clone <url-del-repositorio> pymekit && cd pymekit
pnpm install                # dependencias de todo el monorepo
pnpm supabase:web:start     # Postgres, Auth, Storage y Studio en Docker
pnpm dev                    # aplicacion web en http://localhost:3100
\end{lstlisting}

El primer arranque de Supabase aplica todas las migraciones y el
\textit{seed}. Una vez en marcha, quedan disponibles los servicios de la
tabla~\ref{tab:dg-puertos}.

\begin{table}[htbp]
  \centering
  \small
  \caption{Servicios del entorno local.}
  \label{tab:dg-puertos}
  \begin{tabularx}{\textwidth}{@{}l>{\raggedright\arraybackslash}p{4.4cm}X@{}}
    \toprule
    \textbf{Servicio} & \textbf{Dirección} & \textbf{Uso} \\
    \midrule
    Aplicación web & \texttt{http://localhost:3100} & Web pública, área privada, consola y API del CMS (ADR-009) \\
    API de Supabase & \texttt{http://127.0.0.1:54321} & Auth, Storage y API de datos \\
    PostgreSQL & \texttt{127.0.0.1:54322} & Conexión directa a la base de datos \\
    Supabase Studio & \texttt{http://127.0.0.1:54323} & Herramienta de desarrollo para inspeccionar la base de datos \\
    Correo de pruebas & \texttt{http://127.0.0.1:54324} & Bandeja donde llegan los correos de verificación, invitaciones y enlaces \\
    \bottomrule
  \end{tabularx}
\end{table}

\paragraph{Usuarios de desarrollo.} El \textit{seed} crea, con la contraseña
\texttt{testingpassword}, las cuentas \texttt{test@pymekit.test},
\texttt{owner@pymekit.test}, \texttt{member@pymekit.test},
\texttt{custom@pymekit.test}, \texttt{super-admin@pymekit.test} y
\texttt{cms-staff@pymekit.test} (personal del CMS con el rol «Soporte»). Las dos
últimas tienen un factor TOTP verificado cuyo secreto está en
\texttt{apps/web/supabase/seed.sql}; hay que darlo de alta en una aplicación de
códigos para entrar en la consola. Son credenciales públicas de desarrollo,
válidas solo contra la instancia local (B-03).

\paragraph{Variables de entorno.} Los valores públicos van en
\texttt{apps/web/.env} (compartido) y \texttt{.env.development},
\texttt{.env.test} o \texttt{.env.production}. Los secretos
---\texttt{SUPABASE\_SECRET\_KEY}, \texttt{STRIPE\_SECRET\_KEY},
\texttt{STRIPE\_WEBHOOK\_SECRET}, credenciales del correo--- se ponen en
\texttt{apps/web/.env.local}, que está excluido del repositorio, o en las
variables de la plataforma de despliegue. Nunca se suben al repositorio.

\paragraph{Pagos en local.} Para probar la facturación hay que añadir las
claves de prueba de Stripe (\texttt{VITE\_STRIPE\_PUBLISHABLE\_KEY} y los dos
secretos anteriores) y reenviar los \textit{webhooks} a la aplicación con
\texttt{pnpm stripe:listen}, que publica los eventos en
\texttt{/api/billing/webhook}.

\section{Estructura del monorepo}
\label{sec:dg-estructura}

La tabla~\ref{tab:dg-estructura} resume dónde está cada cosa.

\begin{table}[htbp]
  \centering
  \footnotesize
  \caption{Estructura del monorepo.}
  \label{tab:dg-estructura}
  \begin{tabularx}{\textwidth}{@{}>{\raggedright\arraybackslash}p{5.3cm}X@{}}
    \toprule
    \textbf{Ruta} & \textbf{Contenido} \\
    \midrule
    \texttt{apps/web} & Aplicación TanStack Start (React 19): rutas por fichero en \texttt{src/routes}, \textit{server functions} en ficheros \texttt{*.functions.ts}, configuración en \texttt{src/config} \\
    \texttt{apps/web/src/routes/\_marketing} & Web pública: inicio, precios, blog, contacto, páginas legales \\
    \texttt{apps/web/src/routes/\_authenticated} & Área privada: panel, ajustes, miembros, facturación, creación de equipos \\
    \texttt{apps/web/src/routes/admin} & Consola: inicio y gestión de cuentas (\texttt{/admin}) y pantallas del CMS (\texttt{/admin/cms}) \\
    \texttt{apps/web/src/routes/api} & Rutas de servidor: API del CMS (\texttt{/api/cms/\$}), \textit{webhooks} de pagos y de base de datos, estado de salud \\
    \texttt{apps/web/supabase} & Esquemas declarativos, migraciones, \textit{seed}, plantillas de correo y pruebas pgTAP de toda la base de datos \\
    \texttt{apps/e2e} & Pruebas Playwright de la web y del CMS \\
    \texttt{packages/features/*} & Autenticación, cuentas, equipos, administración y notificaciones \\
    \texttt{packages/billing/*} & Núcleo de facturación, pasarela y proveedor Stripe \\
    \texttt{packages/cms/*} & CMS: servicios Drizzle y rutas Hono (solo servidor), componentes de interfaz, filtros, formateadores, permisos \\
    \texttt{packages/function-middleware} & \textit{Middleware} de las \textit{server functions}: usuario autenticado, miembro de equipo, super-admin, rol mínimo y permiso \\
    \texttt{packages/supabase} & Clientes de Supabase y tipos generados de la base de datos \\
    \texttt{packages/ui}, \texttt{packages/i18n} & Componentes de interfaz y textos (catálogos de mensajes) \\
    \texttt{packages/*} (resto) & Correo, plantillas de correo, analítica, monitorización, OTP, políticas, utilidades compartidas \\
    \texttt{tooling}, \texttt{turbo/generators} & Configuración compartida y generadores de paquetes \\
    \texttt{docs/tfg}, \texttt{scripts/tfg} & Documentación interna del proyecto y scripts de control \\
    \bottomrule
  \end{tabularx}
\end{table}

Los paquetes internos se importan con el \textit{scope} \texttt{@pymekit/*}
(los del CMS, \texttt{@pymekit/cms-*}). Los módulos que solo deben ejecutarse
en el servidor llevan el sufijo \texttt{.server.ts}, y las rutas Hono y los
servicios del CMS nunca se importan desde componentes.

\section{Consideraciones generales sobre el desarrollo}
\label{sec:dg-consideraciones}

\begin{itemize}
  \item \textbf{Server functions.} Se definen con \texttt{createServerFn} en
        ficheros \texttt{*.functions.ts}, con la cadena
        \texttt{.middleware(...).validator(esquemaZod).handler(...)}, y se
        llaman desde el cliente con \texttt{useMutation} y
        \texttt{useServerFn}.
  \item \textbf{Clientes de base de datos.} El cliente de servidor y el del
        navegador aplican RLS. El cliente de administración la ignora: se usa
        lo mínimo posible, con validación manual y un comentario que lo
        justifique.
  \item \textbf{Formularios}: TanStack Form con los campos de
        \texttt{@pymekit/ui/field} y validación con Zod, también en el CMS.
  \item \textbf{Textos}: ningún texto visible ---tampoco \texttt{aria-label},
        \texttt{title}, \texttt{alt} ni \texttt{placeholder}--- se escribe en
        el código; se añade una clave en
        \texttt{packages/i18n/src/messages/es/} (sección~\ref{sec:dg-i18n}).
  \item \textbf{Comentarios} en español didáctico: cabecera en cada fichero
        con lógica, documentación de las funciones exportadas y justificación
        de toda decisión de seguridad. Los puntos que se citan en la memoria
        llevan la etiqueta \texttt{[TFG]} con el requisito.
  \item \textbf{Registro}: \texttt{getLogger()} de
        \texttt{@pymekit/shared/logger}; nunca \texttt{console.log} en
        producción.
  \item \textbf{Git}: ramas \texttt{fase-N/tema}, \textit{commits} en español
        con \textit{Conventional Commits} y el módulo como ámbito.
\end{itemize}

\section{Base de datos}
\label{sec:dg-bd}

\subsection{Esquemas, migraciones y tipos}

El esquema se escribe de forma declarativa en
\texttt{apps/web/supabase/schemas/}, en ficheros numerados cuyo prefijo fija el
orden de dependencias: \texttt{00}--\texttt{18} plataforma SaaS,
\texttt{19} blog y \texttt{20}--\texttt{55} CMS (esquema \texttt{cms}). Las
migraciones se generan a partir de los esquemas, nunca a mano:

\begin{lstlisting}[language=bash]
# 1. Editar o crear el fichero de esquema (schemas/NN-nombre.sql)
# 2. Generar la migracion a partir de la diferencia
pnpm --filter web run supabase:db:diff -f nombre_del_cambio
# 3. Aplicarla y regenerar los tipos TypeScript
pnpm --filter web supabase migrations up
pnpm supabase:web:typegen
# 4. Comprobar que no queda diferencia (no debe mencionar el cambio)
pnpm --filter web run supabase:db:diff
\end{lstlisting}

La herramienta de diferencias no genera los \texttt{revoke} por tabla, los
\texttt{grant} por columna ni la opción \texttt{security\_invoker} de las
vistas (B-35): hay que escribirlos a mano en la migración y comprobarlos con
pgTAP. Las migraciones nunca son destructivas. \texttt{pnpm
supabase:web:reset} recrea la base de datos local desde cero, con las
migraciones y el \textit{seed}, y conviene ejecutarlo antes de fiarse de una
ejecución de las pruebas.

\subsection{Reglas de seguridad}

\begin{itemize}
  \item RLS activada en toda tabla nueva, con políticas basadas en las
        funciones auxiliares existentes (\texttt{has\_role\_on\_account},
        \texttt{has\_permission}, \texttt{is\_mfa\_compliant}).
  \item \texttt{revoke all} a \texttt{anon}, \texttt{authenticated} y
        \texttt{service\_role} tras crear la tabla, y después solo los verbos
        que use la funcionalidad. \texttt{UPDATE} siempre por columnas, sin
        incluir \texttt{id} ni \texttt{account\_id}.
  \item Ninguna función \texttt{security definer} sin controles de acceso
        explícitos.
  \item Todo cambio de políticas, \texttt{grant}, funciones
        \texttt{security definer} o vistas lleva pruebas pgTAP y una revisión
        adversarial de RLS.
\end{itemize}

\subsection{Pruebas pgTAP}

Las pruebas están en \texttt{apps/web/supabase/tests/database/*.test.sql} y se
ejecutan con \texttt{pnpm supabase:web:test}. El fichero
\texttt{00000-pymekit-helpers.sql} define el esquema auxiliar
\texttt{pymekit.*}, con funciones como \texttt{pymekit.authenticate\_as(...)}
para simular a un usuario concreto; las del CMS usan además
\texttt{00001-cms-test-helpers.sql}. Una prueba de aislamiento típica se
autentica como miembro de un equipo e intenta leer o modificar datos de otro.

\section{Añadir una tabla de negocio}
\label{sec:dg-tabla}

Este es el caso más habitual al construir un SaaS sobre PymeKit. Como ejemplo,
una tabla de citas (\texttt{appointments}) que pertenece a una cuenta de
equipo y debe gestionarse desde la consola en un área «Agenda».

\paragraph{1. Esquema con RLS.} Se crea
\texttt{apps/web/supabase/schemas/56-appointments.sql} a partir de la
plantilla de tabla de \texttt{apps/web/supabase/AGENTS.md}:

\begin{lstlisting}[language=SQL]
create table if not exists public.appointments (
  id uuid not null default extensions.uuid_generate_v4(),
  account_id uuid references public.accounts(id) on delete cascade not null,
  customer_name varchar(255) not null,
  starts_at timestamptz not null,
  created_at timestamptz not null default now(),
  updated_at timestamptz not null default now(),
  created_by uuid references auth.users on delete set null,
  updated_by uuid references auth.users on delete set null,
  primary key (id)
);

alter table public.appointments enable row level security;
revoke all on public.appointments from anon, authenticated, service_role;
grant select, insert, update, delete on public.appointments to service_role;
grant select, insert on public.appointments to authenticated;

create policy appointments_read on public.appointments for select
  to authenticated using (public.has_role_on_account(account_id));
create policy appointments_insert on public.appointments for insert
  to authenticated with check (public.has_role_on_account(account_id));
create policy restrict_mfa_appointments on public.appointments
  as restrictive to authenticated using (public.is_mfa_compliant());
\end{lstlisting}

La plantilla completa añade además los \textit{triggers} de marcas de tiempo y
de trazabilidad y un índice por \texttt{account\_id}.

\paragraph{2. Migración, tipos y pruebas.} Se genera la migración con
\texttt{supabase:db:diff -f appointments}, se completan a mano los
\texttt{revoke}, se aplica, se regeneran los tipos y se añade
\texttt{tests/database/appointments.test.sql} comprobando que un miembro de
otro equipo no puede leer ni insertar citas ajenas. Si la tabla está en
\texttt{public}, hay que confirmar que \texttt{anon-surface.test.sql} sigue
pasando: el rol anónimo solo puede leer el blog.

\paragraph{3. Registro en el CMS.} Para que la consola conozca la tabla, se
sincronizan sus metadatos. Puede hacerse desde la interfaz («Ajustes»
$\rightarrow$ «Recursos» $\rightarrow$ «Sincronizar tablas») o, para que
quede en todas las instalaciones, en la propia migración:

\begin{lstlisting}[language=SQL]
select cms.sync_managed_tables('public', 'appointments');

update cms.table_metadata
set display_name = 'Citas',
    ui_config = coalesce(ui_config, '{}'::jsonb)
      || jsonb_build_object('navigation_group', 'Agenda')
where schema_name = 'public' and table_name = 'appointments';
\end{lstlisting}

La clave \texttt{navigation\_group} de \texttt{ui\_config} es el área de la
barra lateral: si no existe, se crea. La sincronización conserva lo ya
configurado, de modo que puede repetirse sin perder los cambios. Los nombres
de las columnas, sus descripciones (la ayuda «?»), su formato y el diseño de
la ficha se ajustan después en «Ajustes» $\rightarrow$ «Recursos», sin código
(el ejemplo de la migración \texttt{20261002140000\_cms\_navigation\_groups.sql}
muestra cómo fijarlos también por migración).

\paragraph{4. Permisos del personal.} El super-administrador ve la tabla
inmediatamente. Para el resto del personal, se crean o amplían roles en
«Ajustes» $\rightarrow$ «Roles y permisos» con permisos de datos sobre
\texttt{public.appointments}. El CMS accede a las tablas mediante funciones
del esquema \texttt{cms} que comprueban esos permisos y registran cada escritura
en la auditoría; las políticas RLS de la tabla siguen protegiendo el acceso de
los clientes desde la aplicación.

\paragraph{5. Interfaz del cliente.} Si los clientes también deben ver las
citas, se añade una ruta en \texttt{apps/web/src/routes/\_authenticated}, una
\textit{server function} con \texttt{teamAccountFunctionMiddleware} y los
textos en el catálogo de mensajes.

\section{Configuración reutilizable}
\label{sec:dg-configuracion}

Adaptar PymeKit a un producto nuevo empieza por la configuración, sin tocar el
núcleo (RNF-01). Los ficheros de \texttt{apps/web/src/config} validan sus
valores con Zod al arrancar, de modo que un error de configuración se detecta
de inmediato (tabla~\ref{tab:dg-configuracion}).

\begin{table}[htbp]
  \centering
  \footnotesize
  \caption{Ficheros de configuración y variables principales.}
  \label{tab:dg-configuracion}
  \begin{tabularx}{\textwidth}{@{}>{\raggedright\arraybackslash}p{4.3cm}X@{}}
    \toprule
    \textbf{Fichero} & \textbf{Qué configura (variables)} \\
    \midrule
    \texttt{app.config.ts} & Nombre, título, descripción, URL, tema y colores (\texttt{VITE\_PRODUCT\_NAME}, \texttt{VITE\_SITE\_TITLE}, \texttt{VITE\_SITE\_DESCRIPTION}, \texttt{VITE\_SITE\_URL}, \texttt{VITE\_DEFAULT\_THEME\_MODE}, \texttt{VITE\_THEME\_COLOR}, \texttt{VITE\_THEME\_COLOR\_DARK}); en producción exige HTTPS \\
    \texttt{account-mode.config.ts} & Modo de cuentas: \texttt{personal-only}, \texttt{organizations-only} o \texttt{hybrid} (\texttt{VITE\_ACCOUNT\_MODE}); es la fuente única de qué superficies existen \\
    \texttt{feature-flags.config.ts} & Interruptores: selector de tema, borrado de cuentas personales y de equipo, creación de equipos, facturación personal y de equipo, notificaciones (\texttt{VITE\_ENABLE\_*}) \\
    \texttt{auth.config.ts} & Métodos de acceso: contraseña, enlace mágico, CAPTCHA (\texttt{VITE\_AUTH\_PASSWORD}, \texttt{VITE\_AUTH\_MAGIC\_LINK}, \texttt{VITE\_CAPTCHA\_SITE\_KEY}) \\
    \texttt{billing.config.ts} & Proveedor (\texttt{VITE\_BILLING\_PROVIDER}), planes, periodicidades, precios e identificadores de precio de la pasarela \\
    \texttt{paths.config.ts} & Rutas de la aplicación (inicio de sesión, panel, ajustes, facturación\ldots) \\
    \texttt{navigation.config.tsx}, \texttt{marketing-navigation.config.tsx} & Menús del área privada y de la web pública \\
    \texttt{tfg.config.ts} & Bloque «Proyecto académico» y recorrido de la demostración en la página de inicio (\texttt{VITE\_SHOW\_TFG\_INFO}); una pyme real lo desactiva con \texttt{false} \\
    \bottomrule
  \end{tabularx}
\end{table}

A esta configuración por ficheros se suma la configuración \emph{como datos}
de la consola ---nombres de tablas y columnas, áreas, formatos, diseño de las
fichas y roles del personal---, que se cambia desde la propia interfaz
(sección~\ref{sec:dg-tabla}). El puerto de la aplicación (3100) se fija en
\texttt{vite.config.ts} y en la variable \texttt{PORT}.

\section{Internacionalización}
\label{sec:dg-i18n}

La interfaz solo ofrece el español, pero todos sus textos salen de catálogos
de mensajes (ADR-021). Para añadir un texto, se crea la clave en inglés en
\texttt{packages/i18n/src/messages/es/<namespace>.json} con el valor en español
y se usa con el componente \texttt{Trans} o con las funciones de mensajes. Las
etiquetas de los valores de los enumerados de la base de datos van en
\texttt{common.enums.<enumerado>.<valor>}.

Añadir un idioma es un punto de extensión documentado en
\texttt{packages/i18n/README.md}. En resumen: copiar
\texttt{src/messages/es/} con el código del idioma nuevo y traducir los
valores; registrarlo en \texttt{src/messages/index.ts} y en
\texttt{supportedLocales} de \texttt{src/config.ts}; crear las traducciones de
los correos en \texttt{packages/email-templates}; elegir el idioma de los
mensajes de validación; volver a mostrar el selector de idioma; y ejecutar
\texttt{pnpm --filter @pymekit/i18n test:unit}, cuya prueba comprueba que el
nuevo catálogo tiene exactamente las mismas claves y marcadores que el
español.

\section{Pruebas}
\label{sec:dg-pruebas}

\begin{lstlisting}[language=bash]
pnpm typecheck              # tipos de todos los paquetes
pnpm lint:fix && pnpm format:fix
pnpm test:unit              # pruebas unitarias (Vitest)
pnpm supabase:web:test      # pruebas de base de datos (pgTAP)

# Pruebas E2E: contra la build de test, no contra el servidor de desarrollo
pnpm --filter web build:test
pnpm --filter web start:test &
pnpm --filter web-e2e test
\end{lstlisting}

Las pruebas E2E se lanzan contra la \textit{build} de test porque el inicio de
sesión automatizado falla contra el servidor de desarrollo (B-13); con
\texttt{pnpm --filter web-e2e test:ui} se abre la interfaz gráfica de Playwright.
Antes de dar por terminado un cambio se completa además la verificación del
proyecto: los controles automáticos del repositorio (\texttt{scripts/tfg/}),
revisión del cambio y, si se tocó la base de
datos, revisión de RLS. La integración continua
(\texttt{.github/workflows/ci.yml}) repite los controles, los tipos, el
\textit{lint}, el formato y las pruebas unitarias en cada \textit{push}, y la
base de datos y las E2E en las \textit{pull requests} hacia \texttt{main}.

\section{Instrucciones para la construcción y el despliegue}
\label{sec:dg-construccion}

\begin{lstlisting}[language=bash]
pnpm build                  # construye todos los paquetes y la web
pnpm --filter web start     # sirve la build de produccion (.output/)
\end{lstlisting}

El servidor de producción lee las variables de \texttt{apps/web/.env} y del
entorno, escucha en el puerto de \texttt{PORT} y expone
\texttt{/api/healthcheck} para las comprobaciones de salud. La base de datos de
un proyecto de Supabase en la nube se actualiza con
\texttt{pnpm --filter web supabase:deploy}, que enlaza el proyecto
(\texttt{SUPABASE\_PROJECT\_REF}) y aplica las migraciones pendientes. El
fichero \texttt{railway.json} define la construcción y el arranque en la
plataforma de despliegue. El procedimiento completo, con la configuración de
Supabase y de Stripe en la nube, se describe en el capítulo de Despliegue del
Sistema (capítulo~\ref{cap:despliegue}).
